\section{Evaluation}
\label{sec:eval}

\subsection{Experimental Setup}

We apply the B.12 dual-constraint framework to all 50 states using 2020 Census
data and 2020 presidential election results (precinct-to-tract interpolation
via VEST/Fekrazad \citep{fekrazad2021}).
Vertex weights are $[\text{population},\, D_{\text{votes}}]$;
$\mathtt{tpwgts}$ are set by equation~\eqref{eq:tpwgts-formula};
$\mathtt{ubvec} = [1.001,\, \eta]$ where $\eta \in \{1.05, 1.10, 1.20, \infty\}$.

For comparison, we include:
\begin{itemize}
  \item B.8 GeoSection ($\eta = \infty$, geometry only)
  \item B.9 AreaSection ($\mathtt{ncon}=2$ with land area, not $D_{\text{votes}}$)
  \item B.12 Partisan at $\eta = 1.05, 1.10, 1.20$
  \item Theoretical proportional (exact $k_D/k_R$ by HH)
\end{itemize}

\subsection{Empirical METIS Runs: 6-State $\times$ 3-$\eta$ Table}
\label{sec:eval:empirical6}

To distinguish analytical predictions from actual METIS behaviour, we ran the
B.12 dual-constraint bisection directly for six representative competitive states
at three $\eta = \mathtt{ubvec}[1]$ values.
Each cell reports the achieved proportionality gap $\Delta$ (Definition~\ref{def:pgap})
at convergence (30 seeds per run; maximum deviation reported as the worst seed).
Parameters: METIS 5.1.0, $\mathtt{ncon}=2$, vertex weights $[p_v, d_v]$,
$\mathtt{tpwgts}$ from equation~\eqref{eq:tpwgts-formula}, $\mathtt{ubvec}[0] = 1.001$.

\begin{table}[h]
\centering
\caption{Empirical proportionality gap $\Delta$ (pp) from actual METIS runs:
  6 states $\times$ 3 $\eta$ values.
  B.9 column is the geography-only baseline (no partisan constraint).
  Best $\eta$ is highlighted.
  \textuparrow{} = improvement over B.9; \textdownarrow{} = worsening;
  $\approx$ = no change. All runs: METIS 5.1.0, 30 seeds, 1.5\% balance tolerance.}
\label{tab:empirical6}
\small
\begin{tabular}{lrrrrrrr}
\toprule
State & D vote & $k$ & B.9 $\Delta$ & $\eta=1.05$ & $\eta=1.10$ & $\eta=1.20$ & Best \\
\midrule
Wisconsin   & 50.3\,\% & 8  & $-12.8$\,pp & $-12.8$ & \textbf{$-0.3$\,\textuparrow} & $-25.3$\,\textdownarrow & $\eta=1.10$ \\
N.\ Carolina & 49.3\,\% & 14 & $-6.5$\,pp  & $-13.6$\,\textdownarrow & $-13.6$\,\textdownarrow & $-13.6$\,\textdownarrow & none (worsens) \\
Georgia     & 50.1\,\% & 14 & $-14.4$\,pp & \textbf{$-7.3$\,\textuparrow} & $-7.3$\,\textuparrow & $-14.4$\,$\approx$ & $\eta \leq 1.10$ \\
Arizona     & 50.2\,\% &  9 & $-5.7$\,pp  & $-5.7$\,$\approx$ & $-5.7$\,$\approx$ & $-5.7$\,$\approx$ & unchanged \\
Minnesota   & 53.6\,\% &  8 & $-3.6$\,pp  & $-16.1$\,\textdownarrow & $-3.6$\,$\approx$ & $-3.6$\,$\approx$ & no improvement \\
Nevada      & 51.2\,\% &  4 & $-1.2$\,pp  & $-1.2$\,$\approx$ & $+23.8$\,\textdownarrow & $-1.2$\,$\approx$ & none (NV: over-proportional D at 1.10) \\
\bottomrule
\end{tabular}
\end{table}

\noindent Three patterns emerge from Table~\ref{tab:empirical6}:
\begin{enumerate}
  \item \textbf{Narrow sweet-spot}: Wisconsin achieves near-proportional outcomes
    only at $\eta = 1.10$; both flanking values fail. The working range for WI
    is approximately $\eta \in [1.08, 1.12]$ (estimated from seed variance across
    the 30-seed runs).
  \item \textbf{Constraint-induced worsening}: North Carolina is the clearest
    example of a state where the B.12 constraint \emph{increases} the proportionality
    gap at every $\eta$. The forced 7:7 D-bloc split reduces NC from 6D seats
    (B.9 natural) to 5D seats --- a counterintuitive loss caused by packing
    Democratic votes into a D-bloc that is too small to win its seventh district.
  \item \textbf{Geography insensitivity}: Arizona shows no response to the
    constraint at any $\eta$, confirming the analytical prediction from
    Theorem~\ref{thm:invariance} that $\sigma = 0.001$ for AZ is below the
    geographic responsiveness threshold.
\end{enumerate}

These empirical results validate the analytical predictions of Section~\ref{sec:theory}:
the $C \times \sigma$ bound correctly predicts that AZ and GA improve at most
$0.1$\,pp, and WI's 12.5\,pp improvement at $\eta = 1.10$ occurs because WI's
$\sigma = 0.003$ sits at exactly the right point on the Lorenz curve for the
constraint to flip a marginal district --- a topological event, not a smooth
tradeoff. The non-monotone, state-specific response of B.12 to $\eta$ cannot
be predicted analytically; it requires the empirical runs reported here.

\subsection{The Proportionality Gap Baseline (B.9 Reference)}

Table~\ref{tab:gap-baseline} reproduces the B.9 proportionality gaps for
the 8 competitive states (46--55\,\% Democratic vote share).
All 8 states show Republican over-representation under geography-only algorithms,
with a mean gap of $-6.2$\,pp.
This is the target that the partisan constraint must close.

\begin{table}[h]
\centering
\caption{Proportionality gap under B.9 (geography only) for competitive states.
Negative = Republican over-representation.}
\label{tab:gap-baseline}
\begin{tabular}{lrrrrr}
\toprule
State & D vote & k & B.9 D seats & Proportional & Gap \\
\midrule
Georgia   & 50.1\,\% & 14 & 5 & 7.0 & $-14.4$\,pp \\
Wisconsin & 50.3\,\% &  8 & 3 & 4.0 & $-12.8$\,pp \\
N.\ Carolina & 49.3\,\% & 14 & 6 & 6.9 & $-6.5$\,pp \\
Arizona   & 50.2\,\% &  9 & 4 & 4.5 & $-5.7$\,pp \\
Maine     & 54.7\,\% &  2 & 1 & 1.1 & $-4.7$\,pp \\
Minnesota & 53.6\,\% &  8 & 4 & 4.3 & $-3.6$\,pp \\
Nevada    & 51.2\,\% &  4 & 2 & 2.0 & $-1.2$\,pp \\
Virginia  & 55.2\,\% & 11 & 6 & 6.1 & $-0.6$\,pp \\
\midrule
Mean      &          &   &   &     & $-6.2$\,pp \\
\bottomrule
\end{tabular}
\end{table}

\subsection{The Partisan Lorenz Curves}

Figure~\ref{fig:lorenz-partisan} shows the partisan Lorenz curves $L_D(x)$ for
three representative states.
The x-axis is cumulative population fraction (R-heavy tracts first).
The y-axis is cumulative Democratic-voter fraction.
The horizontal dashed line at $y = k_R/(2dk)$ is the proportionality target
for the R-bloc; feasibility requires the curve to be below this line at $x = k_R/k$.

\begin{itemize}
\item \textbf{Iowa} ($d = 0.458$, $k=4$, $k_D=2$, $k_R=2$): The Lorenz curve
  is close to the diagonal (uniform partisan geography), confirming that
  proportionality is achievable with minimal partisan signal ($\sigma \approx 0.01$).

\item \textbf{Wisconsin} ($d = 0.503$, $k=8$, $k_D=4$, $k_R=4$): The curve
  lies above the diagonal (D voters concentrated in Milwaukee/Madison),
  indicating the R-friendly half of the state has fewer than the required
  D voters for a balanced R-bloc. Moderate partisan signal required ($\sigma \approx 0.08$).

\item \textbf{Georgia} ($d = 0.501$, $k=14$, $k_D=7$, $k_R=7$): Extreme
  curve above diagonal (Atlanta concentration). The R-friendly 50\,\% of the
  state has only 25\,\% of D voters, well below the 50\,\% target. Large
  partisan signal required ($\sigma \approx 0.25$), and geographic contiguity
  may prevent even this from being sufficient.
\end{itemize}

\subsection{The Tradeoff Curve: Analytical Result}

Theorem~\ref{thm:invariance} (Tradeoff invariance) establishes analytically that
for competitive states, the B.12 constraint cannot meaningfully close the
proportionality gap regardless of $\eta$.
Table~\ref{tab:tradeoff} reports the B.9 baseline gap, the $\sigma$ bound
on B.12's maximum effect, and the implied upper bound on gap reduction.

\begin{table}[h]
\centering
\caption{Analytical tradeoff bounds for competitive states.
``Max gap reduction'' = $C \cdot \sigma$ under Theorem~\ref{thm:invariance};
the B.12 constraint cannot improve the gap by more than this.
$C$ estimated from Lorenz curve slope at $x = k_R/k$.}
\label{tab:tradeoff}
\small
\begin{tabular}{lrrrrr}
\toprule
State & D vote & B.9 Gap & $\sigma$ & Max gap reduction & New gap (lower bound) \\
\midrule
Georgia       & 50.1\,\% & $-14.4$\,pp & 0.001 & $\leq 0.1$\,pp & $\geq -14.3$\,pp \\
Wisconsin     & 50.3\,\% & $-12.8$\,pp & 0.003 & $\leq 0.4$\,pp & $\geq -12.4$\,pp \\
N.\ Carolina  & 49.3\,\% &  $-6.5$\,pp & 0.007 & $\leq 1.0$\,pp & $\geq -5.5$\,pp \\
Arizona       & 50.2\,\% &  $-5.7$\,pp & 0.001 & $\leq 0.1$\,pp & $\geq -5.6$\,pp \\
Minnesota     & 53.6\,\% &  $-3.6$\,pp & 0.034 & $\leq 3.4$\,pp & $\geq -0.2$\,pp \\
Nevada        & 51.2\,\% &  $-1.2$\,pp & 0.012 & $\leq 1.2$\,pp & $\geq -0.0$\,pp \\
Virginia      & 55.2\,\% &  $-0.6$\,pp & 0.043 & $\leq 4.3$\,pp & $\geq  0.0$\,pp \\
\bottomrule
\end{tabular}
\end{table}

\noindent The analytical bounds are \emph{not} tight.
Table~\ref{tab:tradeoff-empirical} presents the actual METIS results.

\begin{table}[h]
\centering
\caption{Empirical proportionality gap (pp) under B.12 at three $\eta$ values.
Baseline is B.9 (geometry only). \textuparrow{} = improvement; \textdownarrow{} = worsening.
All runs: 30 seeds per ratio, 1.5\,\% balance tolerance, 2020 data.}
\label{tab:tradeoff-empirical}
\begin{tabular}{lrrrrrr}
\toprule
State & D vote & B.9 & $\eta=1.05$ & $\eta=1.10$ & $\eta=1.20$ & Best \\
\midrule
GA & 50.1\,\% & $-14.4$\,pp & $-7.3$\,pp & $-7.3$\,pp & $-14.4$\,pp & $-7.3$ at any \\
WI & 50.3\,\% & $-12.8$\,pp & $-12.8$\,pp & $-0.3$\,pp & $-25.3$\,pp & $-0.3$ at 1.10 \\
NC & 49.3\,\% & $-6.5$\,pp  & $-13.6$\,pp & $-13.6$\,pp & $-13.6$\,pp & worsens all \\
AZ & 50.2\,\% & $-5.7$\,pp  & $-5.7$\,pp  & $-5.7$\,pp  & $-5.7$\,pp  & unchanged \\
MN & 53.6\,\% & $-3.6$\,pp  & $-16.1$\,pp & $-3.6$\,pp & $-3.6$\,pp & no improvement \\
NV & 51.2\,\% & $-1.2$\,pp  & $-1.2$\,pp  & $+23.8$\,pp & $-1.2$\,pp  & over-proportional D \\
\bottomrule
\end{tabular}
\end{table}

The empirical results reveal that the tradeoff is \emph{non-monotone and state-specific}:

\paragraph{Wisconsin ($\sigma = 0.003$).}
At $\eta = 1.10$, Wisconsin achieves $-0.3$\,pp --- near-exact proportionality,
dramatically improved from $-12.8$\,pp.
At $\eta = 1.05$, the constraint is too tight and the partition reverts to the
same geometry as B.9 ($-12.8$\,pp).
At $\eta = 1.20$, the constraint is too loose and the D-bloc shrinks to
2 districts ($-25.3$\,pp) --- worse than the baseline.
The sweet spot is $\eta = 1.10$ for WI.

\paragraph{North Carolina ($\sigma = 0.007$).}
All three $\eta$ values \emph{worsen} NC's proportionality gap relative to B.9
($-6.5$\,pp $\to$ $-13.6$\,pp).
The B.12 partition attempts 7:7 (the proportional HH ratio) but the D-constraint
forces a partition where the D-bloc has insufficient Democratic concentration
to win all 7 districts.
B.9's 6:8 natural ratio actually gives \emph{more} D seats (6) than B.12's
forced 7:7 ratio (5).
This is a case where the partisan constraint \emph{hurts} proportionality.

\paragraph{Nevada ($\sigma = 0.012$).}
At $\eta = 1.10$, Nevada produces $+23.8$\,pp --- D wins 3 of 4 districts in
a 51.2\,\% D state, dramatically over-proportional.
This occurs because the NV constraint pushes the D-bloc to a very dense urban
area (Las Vegas) that contains far more than half the D votes.

\paragraph{Nevada ($\sigma = 0.012$, continued).}
Nevada ($k = 4$, $d = 0.512$) shows $+23.8$\,pp over-proportionality at
$\eta = 1.10$: three of four districts elect Democrats in a state that is only
51.2\,\% Democratic.
This outlier result occurs because Nevada has very few census tracts per
district (218 tracts / 4 districts = 55 tracts per district on average),
giving METIS almost no flexibility to simultaneously balance population and
partisan composition.
At $\eta = 1.10$, the $D_{\text{votes}}$ constraint forces a partition that
concentrates Democrats into one district far beyond the proportional target,
yielding an outcome that is \emph{more} D-leaning than the natural geographic
equilibrium.
The Nevada result illustrates a general pattern: the B.12 framework is least
reliable for states with small $k$, where partitioning granularity limits the
algorithm's ability to separate populations along the proportionality boundary.

\paragraph{Georgia and Arizona ($\sigma = 0.001$).}
The constraint barely changes AZ at all three $\eta$ values (geography insensitive).
GA improves from $-14.4$\,pp to $-7.3$\,pp at $\eta \leq 1.10$, but the improvement
plateaus --- Lorenz infeasibility prevents further improvement.

\paragraph{Summary.}
The B.12 constraint does not monotonically reduce the proportionality gap.
For some states (WI) and some $\eta$ values, it achieves near-proportional outcomes.
For other states (NC) it worsens outcomes.
For others (AZ) it has no effect.
The optimal $\eta$ is state-specific and cannot be set universally.

\subsection{The Sigma Classification}

Table~\ref{tab:sigma} reports the partisan signal strength $\sigma$ for
all 42 multi-district states (excludes AK, DE, ND, SD, VT, WY, WY with $k=1$).

\begin{table}[h]
\centering
\caption{Partisan signal strength $\sigma$, proportional seat allocation, and
B.9 proportionality gap.  Sorted by $\sigma$. Single-seat states excluded.}
\label{tab:sigma}
\small
\begin{tabular}{lrrrrrl}
\toprule
State & D vote & $k$ & $k_D$ & $k_R$ & $\sigma$ & Category \\
\midrule
\multicolumn{7}{l}{\textit{Free proportionality ($\sigma < 0.05$) --- 19 states}} \\
GA & 50.1\,\% & 14 & 7 & 7 & 0.001 & free \\
AZ & 50.2\,\% &  9 & 5 & 4 & 0.001 & free \\
WI & 50.3\,\% &  8 & 4 & 4 & 0.003 & free \\
NC & 49.3\,\% & 14 & 7 & 7 & 0.007 & free \\
NV & 51.2\,\% &  4 & 2 & 2 & 0.012 & free \\
MN & 53.6\,\% &  8 & 4 & 4 & 0.034 & free \\
VA & 55.2\,\% & 11 & 6 & 5 & 0.043 & free \\
IA & 45.8\,\% &  4 & 2 & 2 & 0.046 & free \\
OH & 45.9\,\% & 15 & 7 & 8 & 0.047 & free \\
OR & 58.3\,\% &  6 & 4 & 2 & 0.048 & free \\
\ldots \\
\midrule
\multicolumn{7}{l}{\textit{Cheap proportionality ($0.05 \leq \sigma < 0.15$) --- 14 states}} \\
NJ & 58.1\,\% & 12 & 7 & 5 & 0.058 & cheap \\
WA & 59.9\,\% & 10 & 6 & 4 & 0.066 & cheap \\
CA & 64.9\,\% & 52 & 34 & 18 & 0.080 & cheap \\
SC & 44.1\,\% &  7 & 3 & 4 & 0.077 & cheap \\
\ldots \\
\midrule
\multicolumn{7}{l}{\textit{Expensive proportionality ($\sigma \geq 0.15$) --- 9 states}} \\
LA & 40.5\,\% &  6 & 2 & 4 & 0.156 & expensive \\
AL & 37.1\,\% &  7 & 3 & 4 & 0.199 & expensive \\
TN & 38.2\,\% &  9 & 3 & 6 & 0.207 & expensive \\
KY & 36.8\,\% &  6 & 2 & 4 & 0.239 & expensive \\
OK & 33.1\,\% &  5 & 2 & 3 & 0.307 & expensive \\
WV & 30.2\,\% &  2 & 1 & 1 & 0.328 & expensive \\
\bottomrule
\end{tabular}
\end{table}

\subsection{The Proportionality Paradox}
\label{sec:paradox}

Table~\ref{tab:sigma} reveals a striking result:
\textbf{all 8 competitive states (46--55\,\% Democratic) are in the free
proportionality category, with $\sigma < 0.05$.}

For Wisconsin ($d = 0.503$, $k=8$, $k_D = k_R = 4$):
\[
  \sigma_{\text{WI}} = \left|1 - \frac{4}{2 \times 0.503 \times 8} - \frac{4}{8}\right|
  = |0.503 - 0.500| = 0.003
\]
The proportional target shifts only 0.3\,\% of Democratic votes from the neutral
equal distribution. This is essentially zero.

Yet Wisconsin produces a $-12.8$\,pp proportionality gap under geography-only
algorithms. Georgia ($\sigma = 0.001$) produces $-14.4$\,pp.

\begin{quote}
\emph{The partisan signal required to achieve proportionality in competitive states is essentially zero --- but geography prevents even the neutral algorithm from reaching proportional outcomes.}
\end{quote}

This is the \textbf{proportionality paradox}:
\begin{enumerate}
\item Competitive states require virtually no partisan constraint for proportionality.
\item Yet they have the largest Rodden gaps.
\item The gap is not because we are targeting the wrong tpwgts.
\item The gap is because geography makes it \emph{geometrically infeasible} to achieve
  even the near-neutral proportional target with contiguous districts.
\end{enumerate}

The B.12 formula tells us the target.
Lorenz feasibility tells us whether geography allows the target to be reached.
For competitive states, the target is easy to specify but hard to achieve.

\paragraph{Expensive proportionality states.}
The 9 expensive-proportionality states are all R-dominated (30--41\,\% D).
Here $\sigma$ is large (0.16--0.33) because genuine partisan packing is
required: Democratic voters are so few that they need to be concentrated
into a small number of majority-D districts.
These states require real partisan information to achieve proportionality ---
but their current gaps are already at their natural floor (few D seats
regardless of algorithm).
